% !TEX root = main.tex
We consider a $T$-period discrete-time system to model the AI-human pipeline on a platform. Each job has a type in a set $\set{K}$ with $|\set{K}| = K$. In period $t=1,\ldots,T$, at most one new job arrives with probability $\lambda(t) \leq 1$ and we label it as job $t$. The type of job $t$, $\kappa(t)$, is equal to $k \in \set{K}$ with probability $\lambda_k(t) / \lambda(t)$. We use $\Lambda_k(t) = \indic{\kappa(t) = k}$ to denote whether a type-$k$ job arrives in period~$t$. We call $\lambda_k(t)$ the arrival rate of a type-$k$ job in period $t$. For notation convenience, when there is no job arrival, we denote $\kappa(t) = \perp.$

A job $t$ has an inherent cost $C_t \in \mathbb{R}$; 
this quantity being negative means that the platform should accept this job while a positive cost means that the platform should reject this job. If job~$t$ has type~$k$, its cost $C_t$ is independently sampled from an unknown sub-Gaussian distribution $\set{F}_k$ with mean $c_k$. To ease notation, if $\kappa(t) = \perp$, then $C_t \equiv 0$, i.e., the job has zero cost. To distinguish the ability of AI and humans, although the type $\kappa(t)$ is observable, we assume the cost~$C_t$ is
unknown until humans review the job. We use $\expectsub{k}{\cdot}$ to condition on a job type $k$, i.e., $\expectsub{k}{\cdot} = \expect{\cdot \mid \kappa(t) = k}$.

Given the unknown costs, the platform resorts to a AI-human pipeline (Figure~\ref{fig:pipeline}) by making three decisions in any period $t$, \emph{classification}, \emph{admission}, and \emph{scheduling}:
\begin{itemize}
\item \emph{Classification}. For job $t$, the platform first makes a classification decision $Y(t)$ such that the platform accepts the job if $Y(t) = 1$ or rejects it if $Y(t) = -1$. A human may later reverse this classification decision if it is incorrect.
\item \emph{Admission}. The platform can admit job $t$ into the human review system; $A(t) = 1$ if it is admitted and $0$ if not. If $A(t) = 1$, job $t$ is enqueued into an initially empty review queue $\set{Q}$. We define $\set{A}(t) = t$ if $A(t) = 1$ and $\set{A}(t) = \perp$ otherwise.
\item \emph{Scheduling}. At the end of period $t$, the platform selects a job from the review queue $\set{Q}$ for humans to review; we denote this job by $M(t)$. To capture the service capacity, we assume that we have $N(t)$ reviewers in period $t$. Let $\psi_k(t) = \indic{\kappa(M(t)) = k}$ indicate whether the platform schedules a type-$k$ job for review. If $M(t)$ is of type $k$, humans successfully review the job in period $t$ with probability $N(t)\mu_k$, where $\mu_k$ is a known type-specific quantity. We assume that $\max_{t,k} N(t)\mu_k \leq 1$ so that the product is a valid probability. Let $S_k(t) = 1$ and $\set{S}(t) = M(t)$ if the review is successful and $S_k(t) = 0, \set{S}(t) = \perp$ otherwise. If the review is successful, we assume human reviewers observe the exact cost $C_{M(t)}$ and reverse the previous classification decision if wrong. If the review is not successful (e.g., requiring more time to review), the true cost is not observed and the job is placed back to the queue and requires future review. Let $\set{Q}(t)$ be the set of jobs in the review queue at the beginning of period $t$. Then $\set{Q}(t+1) = \set{Q}(t) \cup \{\set{A}(t)\} \setminus \{\set{S}(t)\}$. In addition, the dataset of reviewed jobs at the beginning of period $t$, $\set{D}(t)$, is given by $\set{D}(t) = \{(\set{S}(\tau), C_{\set{S}(\tau)})\}_{\tau < t}$.
\end{itemize}

We note that our model makes no assumption on arrival probabilities $\{\lambda_k(t)\}_{k \in \set{K}, t \leq T}$ and review capacities $\{N(t)\}_{t \leq T}$ (except that $\max_{t,k} N(t)\mu_k \leq 1$). There is no explicit \emph{stabilizability} condition as in the queueing literature (see e.g. \cite{yang2023learning}) since the admission decisions endogenously control the actual arrival rate of the human review queue. However, we do assume that the arrival probabilities and reviewer capacities are exogenous, i.e., this is an oblivious adversary setting.

As a modeling choice, we assume at most one job arrives and at most one service completes per period. This assumption is benign when each period represents a small amount of time, as the probability that more than two jobs arrive at the same time or more than two reviewers finish services is negligible. To ease analysis, we aggregate the service capacity of reviewers in each period (the probability of a successfully reviewed type-$k$ job is $N(t) \mu_k$ for period $t$). That said, our algorithms do not rely on these modeling assumptions.

\subsection{Objective}
The key tension in our setting is between a) purely relying on AI for the classification of a job and b) increasing the congestion of the human review queue by admitting this job. The simplest objective that captures this tension induces the misclassification loss for a job that is either a) not admitted or b) admitted but not reviewed at the end of the time horizon. To simplify exposition and facilitate the comparison to prior work in our empirical study of Section~\ref{sec:content_moderation}, the main body of the paper will focus on the aforementioned objective. That said, our algorithm and analysis extend to other objectives that more finely capture the aforementioned tension (see Remark~\ref{remark:holding_loss}).

Formally, a policy $\pi$ incurs \emph{loss} for any job $t$ that has a wrong classification \emph{at the end} of the horizon and the loss is equal to the absolute value of its cost $|C_t|$. A job has wrong classification if and only if both of the following events happen: (1) this job has an incorrect initial AI classification and (2) this job is not reviewed by a human at the end of the horizon. The latter happens if this job is either not admitted for human review ($A(t) = 0$) or is in the queue in period $T + 1$. Define  the sign function $\sign(x) = \indic{x > 0} - \indic{x \leq 0}.$ The loss of a policy $\pi$ is thus
\begin{equation}\label{eq:policy-loss}
\set{L}^\pi(T) \coloneqq \sum_{t=1}^T |C_t|\underbrace{\mathbbm{1}\Big(Y(t) \neq \sign(C_t)\Big)}_{\text{incorrect AI classification}}\underbrace{\Big(1 - A(t)\indic{t \not \in \set{Q}(T+1)}\Big)}_{\text{not reviewed by humans}}.
\end{equation}

Our goal is to design a \emph{feasible} policy with small loss. We assume that, for each type $k$, its cost distribution $\set{F}_k$ is \emph{unknown} initially and must be learned via the human-reviewed data set $\set{D}(t)$. The platform has no information of $\{\lambda_k(t),N(t)\}_{t \in [T]}$. We assume $\expectsub{k}{|C_t|}$ is bounded by a constant $c_{\max}$ for any $k \in \set{K}$ and that the cost distribution $\set{F}_k$ is sub-Gaussian with a known variance proxy $\sigma_{\max}^2$ such that $\expectsub{k}{\exp(s(C_j - c_k))} \leq \exp(\sigma_{\max}^2 s^2/2)$ for any $s \in \mathbb{R}$ and $k \in \set{K}.$ A policy is \emph{feasible} if its decisions for any period $t$ are only based on the observed sample path $\{\kappa(t'),\set{A}(t'),\set{S}(t')\}_{t' < t} \cup \{\kappa(t)\}$, the dataset $\set{D}(t)$ and the initial information $\sigma_{\max}, c_{\max}$.

\begin{remark}\label{remark:holding_loss}
Our model assumes that all jobs eventually reviewed by humans incur zero loss. However, in practice, the loss of a job misclassified by AI may also depend on its wait time until a human reviews it. For example, in social media platforms, a policy-violating post may get user views before a reviewer removes this post. This loss can be captured by the notion of a \emph{holding cost} in the queueing literature. Our results extend to a setting where the holding cost of a job in any period is bounded by a quantity that vanishes with $T$ (see Appendix~\ref{app:holding-cost}). Subsequent works by \cite{lee2024design, gocmen2025scheduling} design scheduling algorithms minimizing more complex holding costs in the context of content moderation, allowing for posts to have non-vanishing holding cost in a particular period (albeit both \cite{lee2024design, gocmen2025scheduling} assume exogenous classification and admission decisions).
\end{remark}

\subsection{Benchmark}
Directly minimizing the loss is difficult due to the unknown cost distributions. Therefore, we consider a (fluid) benchmark given in \eqref{eq:fluid} that operates with knowledge of the cost distributions. In this fluid benchmark, for every period $t$, there is a mass of $\lambda_k(t)$ jobs of type $k$. The platform admits a mass of $a_k(t)$ jobs to review and leaves a mass of $\lambda_k(t) - a_k(t)$ jobs not reviewed. Admitted jobs receive human reviews immediately. For a non-admitted job $t$ of type $k$, the expected classification loss of accepting this job is $\ell_k^+ \coloneqq \expectsub{k}{(C_t)^+}$ and the loss of rejecting this job is $\ell_k^- \coloneqq \expectsub{k}{(C_t)^-}$, where $x^+ = \max(x,0)$ and  $x^- = -\min(x,0)$.
The loss of a non-admitted job in the benchmark is thus $\ell_k \coloneqq \min(\ell_k^+, \ell_k^-)$ (assuming a known cost distribution), by rejecting a job if $c_k > 0$ or accepting it if $c_k \leq 0$ (this is because $\ell_k^+ \leq \ell_k^-$ if and only if $\ell_k^+ - \ell_k^- = c_k \leq 0$). The expected loss of the benchmark is then $\sum_{t = 1}^T \sum_{k \in \set{K}} \ell_k(\lambda_k(t) - a_k(t))$. To capture the human capacity constraint, we require the amount of admitted jobs to be bounded by the available capacity \emph{in every period}. Specifically, the platform decides the probability $\nu_k(t)$ of scheduling a type-$k$ job for review in period $t$ (recall that at most one job is scheduled to review per period). Since the review of a  scheduled type-$k$ job is successful with probability $\mu_k N(t)$, the mass (throughput) of reviewed type-$k$ jobs for period $t$ is $\mu_k N(t) \nu_k(t)$. Our capacity constraint requires that $a_k(t) \leq \mu_k N(t)\nu_k(t)$ for every type $k$ and period~$t$.

\begin{equation}\label{eq:fluid}
\tag{Fluid}
\begin{aligned}
\mathcal{L}^\star(T) = &\min_{\{a_k(t),\nu_k(t)\}_{k \in [K], t\in [T]}}\sum_{t=1}^T \sum_{k \in \set{K}} \ell_k(\lambda_k(t) - a_k(t)),\text{s.t.} \\
\quad& \bolds{a_k(t) \leq \mu_k N(t)\nu_k(t),\forall k \in \set{K}, t \in [T]} \\
&a_k(t) \leq \lambda_k(t),~\nu_k(t) \geq 0,~\forall k \in \set{K}, t\in [T] \\
& \sum_{k\in \set{K}} \nu_k(t) \leq 1,~\forall t \in [T].
\end{aligned}
\end{equation}
To measure the performance of a policy, we invoke the notion of \emph{regret} that compares its expected loss $\expect{\set{L}^\pi(T)}$ to the loss of the benchmark $\set{L}^\star(T)$.
\begin{equation}\label{eq:regre}
 \reg^{\pi}(T) = (\expect{\set{L}^{\pi}(T)} - \set{L}^{\star}(T))^+.
 \end{equation}
Our goal is to find an \emph{asymptotically optimal} policy where the regret of a policy is sublinear in $T$, i.e., $\lim_{T \to \infty} \reg^{\pi}(T) / T = 0.$

\begin{remark}\label{remark:w-fluid}
In a setting with stationary arrivals and capacity, \eqref{eq:fluid} is a natural benchmark because it lower bounds the loss of any feasible policy. When arrivals or capacity are non-stationary, \eqref{eq:fluid} is weak as it does not allow a job to be reviewed by later available capacity. For ease of exposition, the main body focuses on deriving performance guarantees based on this fluid benchmark. In Appendix~\ref{app:w-fluid}, we show how our guarantees extend to a series of relaxed benchmarks called $w-$fluid benchmarks motivated by non-stationary queueing systems \citep{borodin2001adversarial, liang2018minimizing, yang2023learning, nguyen2023learning}, where the capacity constraint is satisfied over consecutive windows of periods instead of for each period. 
\end{remark}

\paragraph{Notation.} For ease of exposition when stating our results, we use $\lesssim$ to include dependence only on the number of types $K$, and the time horizon $T$ with other parameters treated as constants. Note that we use $\perp$ to denote an empty element, so a set $\{\perp\}$ should be interpreted as an empty set $\emptyset$. We also follow the convention that $\frac{a}{0} = +\infty$ for any positive $a$. 
For a $d-$dimensional positive semi-definite (PSD) matrix  $\bolds{V}$, we define its corresponding vector norm by $\|\btheta\|_{\bolds{V}} = \sqrt{\btheta^{\trans}\bolds{V}\btheta}$ for any $\btheta \in \mathbb{R}^d$. We denote $\bolds{I}$ as the identity matrix with a suitable dimension, $\mathrm{det}(\bV)$ as the determinant of matrix $\bV$, and $\lambda_{\min}(\bV)$ as the minimum eigenvalue of a PSD matrix $\bV$.